% =============================================================================
% Kapitel 3: Methodik
% =============================================================================

Dieses Kapitel beschreibt die methodische Vorgehensweise der vorliegenden Arbeit.
Zunächst wird der übergeordnete Forschungsrahmen dargestellt, anschließend die systematische Anforderungserhebung erläutert und abschließend die Evaluationsmethodik definiert.

\subsection{Forschungsrahmen: Design Science Research}\label{sec:dsr}

Die Arbeit folgt dem Design-Science-Research-Paradigma (DSR) nach Hevner et al.~\cite{hevner2004designscience}.
DSR zielt darauf ab, durch die Konstruktion und Evaluation von IT-Artefakten sowohl praktische als auch wissenschaftliche Erkenntnisse zu gewinnen.
Im Gegensatz zu rein verhaltensorientierten Ansätzen, die beobachten und erklären, steht bei DSR das \emph{Gestalten} im Vordergrund: Ein Artefakt wird entworfen, implementiert und anhand definierter Kriterien evaluiert.

Das zentrale Artefakt dieser Arbeit ist ein hybrider Analyseansatz, der regelbasierte und Embedding-basierte Textanalyse kombiniert, um Kompetenzformulierungen in deutschen Modulhandbüchern zu klassifizieren, taxonomisch einzuordnen und Verbesserungsvorschläge zu generieren.
Der Ansatz wird als prototypische Weiterentwicklung des bestehenden Systems von Loth und Konert~\cite{loth2022kompetenzeditor} realisiert.

Hevner et al. definieren sieben Richtlinien für Design-Science-Forschung, die im Folgenden auf diese Arbeit bezogen werden:

\begin{enumerate}
  \item \textbf{Design as an Artifact:} Das Artefakt ist der hybride Kompetenzeditor -- ein funktionsfähiges System, das die drei Analyseansätze (regelbasiert, Embedding-basiert, hybrid) implementiert und vergleichend evaluierbar macht.
  \item \textbf{Problem Relevance:} Die Qualitätssicherung von Kompetenzformulierungen ist ein reales Problem der Hochschullehre. Bestehende Werkzeuge nutzen entweder statische Verblisten oder rein datengetriebene Ansätze; eine systematische Kombination beider Paradigmen für den deutschsprachigen Raum fehlt (vgl. Kap.~\ref{sec:grundlagen}).
  \item \textbf{Design Evaluation:} Die Evaluation erfolgt quantitativ auf einem Gold-Standard-Datensatz mit den Metriken Precision, Recall, F1-Score und Precision@k (vgl. Kap.~\ref{sec:eval-methodik}).
  \item \textbf{Research Contributions:} Die Beiträge umfassen (a)~den hybriden Cascading-Ansatz als Entwurfsmuster, (b)~den annotierten Gold-Standard-Datensatz deutschsprachiger Kompetenzformulierungen und (c)~die vergleichende Evaluation dreier Ansätze.
  \item \textbf{Research Rigor:} Die Evaluation nutzt etablierte Klassifikationsmetriken, stratifizierte Datensplits und eine systematische Anforderungserhebung nach Sommerville~\cite{sommerville2016software}.
  \item \textbf{Design as a Search Process:} Der iterative Charakter zeigt sich in der schrittweisen Entwicklung: Zunächst wird der bestehende Prototyp analysiert, dann Anforderungen erhoben, darauf aufbauend die Architektur entworfen und schließlich implementiert und evaluiert.
  \item \textbf{Communication of Research:} Die Ergebnisse werden in dieser Thesis dokumentiert und der Quellcode zur Nachvollziehbarkeit bereitgestellt.
\end{enumerate}

\subsection{Anforderungserhebung}\label{sec:anforderungen}

Die Anforderungserhebung folgt einem systematischen Vorgehen, das vier Quellenkategorien kombiniert.
Sommerville~\cite{sommerville2016software} beschreibt Requirements Engineering als iterativen Prozess aus Ermittlung, Analyse, Spezifikation und Validierung.
Im Kontext dieser Arbeit wird die Ermittlung durch die Kombination mehrerer Perspektiven operationalisiert.

\subsubsection{Quellen der Anforderungen}\label{sec:anf-quellen}

Die Anforderungen werden aus vier Quellenkategorien abgeleitet:

\paragraph{Prototyp-Schwächen (S1--S8).}
Die systematische Analyse des bestehenden Prototyps (vgl. Kap.~\ref{sec:vorarbeit}) identifiziert acht konkrete Schwächen: exakter Stringvergleich statt semantischer Analyse~(S1), geschlossene Verbliste~(S2), fehlende Kontextsensitivität~(S3), keine Klassifikation von Satztypen~(S4), statische Empfehlungen~(S5), fehlende Set-Analyse~(S6), simplifizierter Score~(S7) und fehlendes Konfidenzmaß~(S8).
Jede Schwäche motiviert mindestens eine Anforderung.

\paragraph{Betreuer-Feedback (K1--K5).}
Rückmeldungen von Prof.\ Konert zum Exposé definieren zusätzliche Erwartungen: Precision@k als Evaluationsmetrik für Empfehlungen~(K1), Set-Analyse auf Modulebene~(K2), Konfidenzmaße für alle Zuordnungen~(K3), Hybrid als dritte Vergleichsvariante~(K4) und LLM-gestützte Vorannotation des Gold-Standards~(K5).

\paragraph{Literatur.}
Die Taxonomieforschung (Krathwohl~\cite{krathwohl2002revision}, Stanny~\cite{stanny2016reevaluating}, Cursio~\cite{cursio2013leitfaden}) zeigt, dass Verben mehrdeutig sind und Kontextanalyse erfordern.
Die Embedding-Literatur (Reimers und Gurevych~\cite{reimers2019sentencebert}) liefert die technische Grundlage.
Chiticariu et al.~\cite{chiticariu2013rulebased} belegen, dass regelbasierte Ansätze in kommerziellen Systemen dominant sind, wenn Transparenz und Nachvollziehbarkeit gefordert werden.

\paragraph{Praxisanforderungen.}
HRK nexus~\cite{hrknexus2015lernergebnisse} und Kopf et al.~\cite{kopf2010kompetenzen} definieren Qualitätskriterien für Kompetenzformulierungen, die als fachliche Validierungsgrundlage dienen.

\subsubsection{Anforderungskatalog}\label{sec:anf-katalog}

Aus den genannten Quellen werden sechs funktionale Anforderungen (FA1--FA6) und fünf nicht-funktionale Anforderungen (NFA1--NFA5) abgeleitet.
Tab.~\ref{tab:anforderungen} gibt eine Übersicht mit Priorisierung nach dem MoSCoW-Schema.

\begin{table}[H]
\caption{Anforderungskatalog mit MoSCoW-Priorisierung\label{tab:anforderungen}}
\centering
\small
\begin{tabular}{lllll}
\hline
\textbf{ID} & \textbf{Anforderung} & \textbf{Typ} & \textbf{Prio} & \textbf{Quelle} \\
\hline
FA1 & Klassifikation (K vs.\ Inhalt) & funktional & Must & S4, HRK \\
FA2 & Taxonomie-Zuordnung (Stufe 1--6) & funktional & Must & S1--S3, K4 \\
FA3 & Ähnlichkeitsbasierte Empfehlungen & funktional & Should & S5, K1 \\
FA4 & Set-Analyse auf Modulebene & funktional & Should & S6, K2 \\
FA5 & Echtzeit-Feedback & funktional & Must & P-Stärke \\
FA6 & Persistenz & funktional & Should & P-Stärke \\
\hline
NFA1 & Konfidenzmaße pro Zuordnung & nicht-funkt. & Must & S8, K3 \\
NFA2 & Transparenz der Analyse & nicht-funkt. & Should & Chiticariu \\
NFA3 & Deutsche Sprache & nicht-funkt. & Must & Kontext \\
NFA4 & Technologie-Migration (Nuxt 4) & nicht-funkt. & Must & S-EOL \\
NFA5 & Antwortzeit $< 2$\,s & nicht-funkt. & Must & Usability \\
\hline
\end{tabular}
\end{table}

Die sieben Must-have-Anforderungen (FA1, FA2, FA5, NFA1, NFA3, NFA4, NFA5) definieren den Minimalumfang des Systems.
Die Should-have-Anforderungen (FA3, FA4, FA6, NFA2) erweitern den Funktionsumfang und werden im Rahmen der verfügbaren Zeit umgesetzt.

\subsection{Evaluationsmethodik}\label{sec:eval-methodik}

Die Evaluation vergleicht drei Analyseansätze -- regelbasiert, Embedding-basiert und hybrid -- auf einem gemeinsamen Gold-Standard-Datensatz.
Dieser Abschnitt beschreibt die Konstruktion des Datensatzes und die verwendeten Metriken.

\subsubsection{Gold-Standard-Datensatz}\label{sec:gold-standard}

Die Evaluationsgrundlage bildet ein Datensatz aus über 5\,600~Sätzen, die automatisiert aus Modulhandbüchern von vier deutschen Hochschulen und vier Fachbereichen extrahiert wurden.
Die Datenquellen umfassen die Hochschule Fulda (Angewandte Informatik), die TU Darmstadt (Informatik), die TH Mittelhessen (Maschinenbau, Wirtschaftsingenieurwesen) und die Universität Kassel (Soziale Arbeit).
Die Diversität der Hochschulen und Fachbereiche gewährleistet, dass die Evaluation nicht auf eine einzelne Institution oder Domäne beschränkt ist.

Die Annotation erfolgt in einem zweistufigen Verfahren:
Zunächst werden alle Sätze durch ein Large Language Model vorannotiert (LLM-gestützte Vorannotation nach K5), anschließend wird eine stratifizierte Stichprobe manuell validiert.
Jeder Satz erhält drei Annotationen: den \emph{Satztyp} (Kompetenzformulierung, Inhaltsbeschreibung oder Sonstiges), für Kompetenzformulierungen die \emph{Taxonomiestufe} nach Anderson und Krathwohl~\cite{anderson2001taxonomy} (Stufen 1--6, wobei Mehrfachzuordnungen möglich sind) sowie eine Bewertung der \emph{Formulierungsqualität} (gut, akzeptabel, schlecht).
Die Mehrfachzuordnung bei der Taxonomie trägt dem Umstand Rechnung, dass Sätze mit mehreren Handlungsverben verschiedenen Stufen angehören können.

Der Datensatz wird stratifiziert nach Hochschule und Typ in ein Trainings-Set (80\,\%) und ein Test-Set (20\,\%) aufgeteilt (Seed~42), um reproduzierbare Ergebnisse zu gewährleisten.

\subsubsection{Metriken}\label{sec:metriken}

Für die Klassifikation (FA1) und Taxonomie-Zuordnung (FA2) werden Precision, Recall und F1-Score berechnet.
Da die Klassenverteilung unbalanciert ist (94,6\,\% Kompetenzformulierungen), wird der \emph{gewichtete} F1-Score (\emph{weighted F1}) als primäre Metrik verwendet, der die Metriken jeder Klasse proportional zu ihrer Häufigkeit gewichtet.
Ergänzend werden die klassenweisen F1-Werte berichtet, um die Leistung auf unterrepräsentierten Klassen (insbesondere Stufe~4 \emph{Analysieren} und Stufe~5 \emph{Bewerten}) transparent zu machen.

Für das Empfehlungssystem (FA3) wird Precision@k verwendet: Unter den Top-$k$ vorgeschlagenen Referenzformulierungen wird der Anteil relevanter Ergebnisse gemessen.
Die Relevanz wird durch Übereinstimmung der Taxonomiestufe operationalisiert.

Die Set-Analyse (FA4) wird qualitativ bewertet, da keine quantitative Ground Truth auf Modulebene vorliegt.
